\part{Comprendre}

% ─────────────────────────────────────────────────────────────────
\chapter{Coordonnées célestes}
\label{ch:coords}

\section{La sphère céleste}

Vue du sol, la voûte étoilée se comporte comme une sphère qui tourne d'est en
ouest autour d'un axe passant par les pôles célestes. Les astres sont à des
distances très différentes, mais il suffit de leur position sur cette sphère
fictive pour les repérer et les suivre.

La Terre boucle un tour sur elle-même en \textbf{23~h 56~min 4~s} --- le
\emph{jour sidéral}, plus court que les 24~heures du jour solaire. C'est à
cette cadence-là, et non à celle de l'horloge, qu'une monture doit tourner.

\section{Ascension droite}

L'\textbf{ascension droite} (AD, \emph{right ascension} en anglais) est
l'équivalent céleste de la longitude. Elle se compte en \textbf{heures} et non
en degrés, de 0~h à 24~h, depuis le point vernal.

\begin{formula}
\[
1\,\text{h} = 15\degre{} \qquad 1\,\text{min} = 15' \qquad 1\,\text{s} = 15''
\]
Ainsi AD = 02~h 44~min 08~s vaut $2\times15\degre{} + 44\times0{,}25\degre{} +
8\times\frac{15}{3600}\degre{} = 41{,}03\degre$.
\end{formula}

\begin{info}
L'AD est la coordonnée critique pour le suivi : une erreur d'une seconde de
temps vaut 15 secondes d'arc sur le ciel à l'équateur céleste. Sur un
échantillonnage de 1,6\arcsec/pixel, cela fait une dizaine de pixels de
traînée.
\end{info}

\section{Déclinaison}

La \textbf{déclinaison} (DÉC) est l'équivalent de la latitude : la distance
angulaire au nord ou au sud de l'équateur céleste, de $-90\degre$ à $+90\degre$. Sur une
monture correctement alignée, l'axe de déclinaison ne bouge pas pendant le
suivi.

\section{Angle horaire}

L'\textbf{angle horaire} mesure la distance d'un astre au méridien local. Il
augmente sans cesse avec la rotation de la Terre.

\begin{formula}
\[
\text{HA} = \text{temps sidéral local} - \text{AD}
\]
Une tolérance en AD s'exprime en secondes d'arc ou en secondes de temps. La
conversion dépend de la déclinaison :
\[
\text{secondes de temps} = \frac{\text{secondes d'arc}}{15 \cos(\text{DÉC})}
\]
\end{formula}

\begin{tip}
Près du pôle céleste, une erreur en AD déplace très peu l'étoile sur le
capteur : le cosinus de la déclinaison l'atténue. Les objets circumpolaires
pardonnent donc davantage les défauts de suivi en AD --- mais pas ceux en
déclinaison.
\end{tip}

% ─────────────────────────────────────────────────────────────────
\chapter{Montures équatoriales}
\label{ch:mounts}

\section{Principe}

Une \textbf{monture équatoriale} annule la rotation terrestre en faisant
tourner le tube autour d'un axe parallèle à celui de la Terre : l'axe polaire.

\begin{itemize}
  \item \textbf{Axe polaire (AD)} : tourne à la vitesse sidérale.
  \item \textbf{Axe de déclinaison} : immobile pendant le suivi, en théorie.
\end{itemize}

\begin{formula}
\[
\omega = \frac{1\,296\,000''}{86\,164{,}09\,\text{s}} \approx
15{,}041\,\text{secondes d'arc par seconde}
\]
\end{formula}

Tout écart à cette vitesse, ou tout défaut d'alignement de l'axe polaire, se
traduit par une dérive que l'image enregistre.

\section{Les montures 10Micron}

Les 10Micron (GM1000 HPS et au-delà) portent des \textbf{encodeurs absolus} :
elles connaissent leur position même après une coupure, là où un encodeur
relatif la perd. S'y ajoutent un entraînement par vis sans fin, une correction
d'erreur périodique et surtout un \textbf{modèle de pointage} interne bâti sur
plusieurs dizaines d'étoiles.

C'est ce modèle qui rend possible l'imagerie \textbf{non guidée} : la monture
ne se contente pas de tourner à la vitesse sidérale, elle applique des taux
corrigés qui absorbent la flexion, la réfraction et le défaut résiduel
d'alignement polaire.

\begin{warning}
Conséquence directe, et souvent mal comprise : sur une monture à modèle,
\textbf{l'erreur d'alignement polaire est dans le modèle} et déjà compensée par
les taux de suivi. Une dérive résiduelle en déclinaison n'appelle donc pas un
coup de vis d'azimut, mais des \textbf{points de modèle supplémentaires près de
la cible}.
\end{warning}

Ces montures parlent un \textbf{LX200 étendu} : au-delà des commandes standard,
elles exposent les positions brutes d'axes, les taux appliqués, l'état du
modèle et les données d'environnement. MountMonitor s'appuie sur ces
extensions quand elles sont disponibles.

% ─────────────────────────────────────────────────────────────────
\chapter{Ce qui dégrade le suivi}
\label{ch:errors}

Aucune monture ne suit parfaitement. Les causes se superposent, et chacune
laisse une signature différente dans l'enregistrement.

\section{Erreur périodique}

La vis sans fin n'est jamais parfaitement ronde ni parfaitement centrée. Son
défaut se répète à chaque tour et produit une oscillation de période fixe ---
quelques minutes le plus souvent. C'est précisément ce qu'une analyse de
Fourier sait extraire (chapitre~\ref{ch:fft}).

\section{Alignement polaire}

Si l'axe polaire ne vise pas exactement le pôle céleste, l'astre dérive
lentement en déclinaison, d'autant plus vite que l'on s'éloigne du méridien.
Sur une monture sans modèle, la réponse est mécanique. Sur une monture à
modèle, voir l'avertissement du chapitre précédent.

\section{Réfraction atmosphérique}

L'atmosphère relève les astres d'autant plus qu'ils sont bas : environ
$34'$ à l'horizon, $1'$ à $45\degre$ de hauteur. La réfraction varie donc pendant la
nuit, et sa dérive se confond facilement avec un défaut d'alignement.

\section{Perturbations mécaniques et vibrations}

Vent, câbles qui tirent, sol qui bouge, passage du méridien : autant de
secousses ponctuelles qui n'ont rien de périodique. Elles se voient comme des
pics isolés, et c'est ce qui permet de les distinguer du reste.

\section{Mouvements commandés}

Ceux-là ne sont pas des erreurs du tout. Entre deux poses, le séquenceur
déplace délibérément la monture --- \textbf{dither} pour décorréler le bruit,
\textbf{recentrage} quand la cible a glissé --- et la monture \emph{reste} à sa
nouvelle position. Les confondre avec de l'erreur de suivi est la faute la plus
coûteuse qu'on puisse commettre en lisant un enregistrement, et le
chapitre~\ref{ch:stats} y est consacré.

% ─────────────────────────────────────────────────────────────────
\chapter{Lire les chiffres}
\label{ch:stats}

\begin{warning}
Ce chapitre est le plus important du manuel. Un enregistrement de suivi se lit
de travers avec une facilité déconcertante, et le chiffre qu'on en tire est
alors faux d'un facteur plusieurs centaines.
\end{warning}

\section{Le signal est un escalier}

Voici ce qu'enregistre réellement une nuit d'imagerie non guidée : la monture
tient sa position à une fraction de seconde d'arc, puis le séquenceur la
déplace d'un dither, et elle tient sa nouvelle position tout aussi bien. Le
tracé n'est pas une ligne bruitée autour d'une moyenne : c'est un
\textbf{escalier}.

Un écart-type calculé sur toute la nuit mesure alors la hauteur des marches,
pas la qualité du suivi. Et retirer une droite --- ce que fait un
« détrendage » classique --- ne retire pas un escalier.

\begin{info}
Sur une session réelle de huit heures et demie, l'écart quadratique moyen brut
donnait \textbf{86,3\arcsec}, un verdict désastreux. La monture tenait en
réalité chaque position à \textbf{0,18\arcsec}. Le rapport entre les deux
chiffres est de 480.
\end{info}

\section{Le jitter}

Le \textbf{jitter} est la dispersion que montre la monture \emph{pendant}
qu'elle tient une position. C'est le seul chiffre qui étale une étoile sur le
capteur, et c'est celui sur lequel MountMonitor fonde sa note.

Il se mesure entre deux repositionnements, jamais autour d'une droite. Le
programme découpe l'enregistrement aux sauts, écarte deux échantillons après
chacun --- le mouvement lui-même n'est pas du suivi --- et prend la médiane des
dispersions par palier.

\begin{tip}
Cette façon de faire a une propriété qui vaut d'être notée : une séquence qui
dithère à chaque pose ne sort pas moins bonne qu'une séquence qui n'en fait
jamais. Le chiffre décrit la monture, pas le séquenceur.
\end{tip}

\section{Les mouvements commandés}

Ils sont comptés et classés à part :

\begin{center}
\rowcolors{2}{rowalt}{white}
\begin{tabularx}{\textwidth}{lX}
\toprule
\textbf{Classe} & \textbf{Ce que c'est} \\
\midrule
Dither & Petit déplacement entre deux poses, pour décorréler le bruit. \\
Recentrage & Déplacement plus ample, qui ramène la cible où elle doit être. \\
Non stabilisé & La dispersion qui suit reste plusieurs fois la normale. C'est
celui-là qui mérite un coup d'œil. \\
\bottomrule
\end{tabularx}
\end{center}

\begin{warning}
Un franchissement de seuil n'est pas un mouvement. Un recentrage de 70\arcsec{}
met plusieurs échantillons à s'accomplir et déclenche le seuil à chacun : une
nuit comptait ainsi 98 franchissements pour 53 mouvements réels, avec une
amplitude médiane dix fois trop petite. MountMonitor regroupe les
franchissements séparés de moins de dix secondes.
\end{warning}

\section{La dérive lente}

Elle est donnée par heure et par pose. Sur une monture non guidée, chaque
dither l'annule : ce qui compte n'est donc pas la dérive sur la nuit, mais le
déplacement qu'elle produit \emph{pendant une pose}.

\begin{formula}
\[
\text{déplacement sur une pose} = \text{dérive (\arcsec/h)} \times
\frac{\text{durée de pose (s)}}{3600}
\]
Une dérive de $-16{,}8\arcsec$/h déplace donc l'étoile de $0{,}84\arcsec$ sur
une pose de 180~s.
\end{formula}

\section{L'écart-type glissant}

Calculé en direct sur une fenêtre de 60 à 900 secondes, il sert à voir
\emph{quand} quelque chose s'est produit, pas à noter la nuit. Après un
changement de cible, une coupure d'acquisition ou un repositionnement, sa
fenêtre chevauche le mouvement et le décrit lui, pas le suivi : ces valeurs-là
sont écartées du rapport.

\section{La note}

Elle est fondée sur le jitter combiné, avec des seuils adaptés au type de
monture. Une 10Micron non guidée est jugée sur des seuils plus larges qu'une
monture guidée, parce que ce ne sont pas les mêmes ordres de grandeur.

\begin{center}
\rowcolors{2}{rowalt}{white}
\begin{tabularx}{\textwidth}{lXX}
\toprule
\textbf{Note} & \textbf{Monture guidée} & \textbf{Non guidée de précision} \\
\midrule
Excellent & $< 0{,}5\arcsec$ & $< 2{,}0\arcsec$ \\
Bon       & $< 1{,}5\arcsec$ & $< 4{,}0\arcsec$ \\
Correct   & $< 3{,}0\arcsec$ & $< 8{,}0\arcsec$ \\
\bottomrule
\end{tabularx}
\end{center}

\begin{tip}
Rapportez toujours la note à votre échantillonnage. À 1,6\arcsec/pixel, un
jitter de 0,5\arcsec{} fait trembler l'étoile d'un tiers de pixel : invisible
sur l'image finale.
\end{tip}

% ─────────────────────────────────────────────────────────────────
\chapter{L'analyse de Fourier}
\label{ch:fft}

\section{À quoi elle sert}

Une transformée de Fourier décompose un signal en somme de sinusoïdes. Sur un
enregistrement de suivi, elle fait ressortir ce qui \emph{se répète} --- et une
erreur périodique de vis sans fin se répète, par construction.

MountMonitor cherche dans la plage de \textbf{2 à 15 minutes}, celle où tombent
les périodes de vis des montures courantes.

\section{Deux précautions}

\begin{itemize}
  \item La transformée porte sur les déviations \textbf{détrendées}. Une dérive
  linéaire n'est pas périodique, mais sur une fenêtre finie elle produit une
  grosse composante basse fréquence qui masquerait ce qu'on cherche.
  \item La \textbf{limite de Nyquist} impose d'échantillonner au moins deux
  fois par période. À 2~Hz, les périodes inférieures à la seconde sont hors de
  portée --- ce qui n'est pas gênant, aucune erreur mécanique n'y loge.
\end{itemize}

\section{Lire un spectre}

Un pic net à une période stable désigne une erreur périodique. Un spectre plat
signifie qu'il n'y en a pas de décelable --- le cas des 10Micron après
correction. Des pics larges et mouvants sont du bruit, du vent le plus souvent.

% ═══════════════════════════════════════════════════════════════════
\part{Utiliser}

% ─────────────────────────────────────────────────────────────────
\chapter{Installation}
\label{ch:install}

\section{Systèmes pris en charge}

Les paquets publiés embarquent leur propre Python et leur propre Qt : rien
d'autre n'est à installer. En revanche ils exigent un système assez récent,
parce que les binaires Qt qu'ils contiennent l'exigent.

\begin{center}
\rowcolors{2}{rowalt}{white}
\begin{tabularx}{\textwidth}{lXl}
\toprule
\textbf{Système} & \textbf{Minimum} & \textbf{Architecture} \\
\midrule
Windows & Windows 10 (1809) & x64 \\
macOS & macOS 11 Big Sur & Apple Silicon \emph{et} Intel \\
Linux & glibc 2.28 --- Debian 11, Ubuntu 20.04, RHEL 8 & x86\_64 \\
Linux ARM & glibc 2.39 --- Ubuntu 24.04 & arm64 \\
\bottomrule
\end{tabularx}
\end{center}

\begin{warning}
Choisissez le paquet macOS qui correspond à votre machine :
\texttt{-macos-arm64.dmg} pour un Mac Apple Silicon, \texttt{-macos-x86\_64.dmg}
pour un Mac Intel. Le Python embarqué est propre à une architecture même si Qt,
lui, est universel : le mauvais paquet ne démarrera pas du tout.
\end{warning}

\section{Poser le programme}

Sur Windows, l'installeur s'occupe de tout. Sur macOS, glissez l'application
dans le dossier Applications. Sur Linux, \texttt{installer.sh} l'installe dans
le profil de l'utilisateur, sans \texttt{sudo} ni gestionnaire de paquets.

Les mises à jour sont ensuite vérifiées discrètement à chaque démarrage et
proposées quand il y en a.

\section{Premier essai sans monture}

Le mode simulation permet de tout parcourir sans matériel :

\begin{lstlisting}[language=bash]
MountMonitor --sim-all
\end{lstlisting}

% ─────────────────────────────────────────────────────────────────
\chapter{Réglages}
\label{ch:config}

\section{Connexion}

\begin{center}
\rowcolors{2}{rowalt}{white}
\begin{tabularx}{\textwidth}{lX}
\toprule
\textbf{Protocole} & \textbf{Quand l'employer} \\
\midrule
LX200 TCP/IP & 10Micron en réseau. Port 3492 par défaut. \\
LX200 série & Liaison RS-232, 9600 8N1 d'origine sur les 10Micron. \\
ASCOM & Windows uniquement, via le pilote du constructeur. \\
\bottomrule
\end{tabularx}
\end{center}

\section{Le site d'observation}

\begin{warning}
Renseignez \textbf{latitude, longitude et altitude} dans les préférences. La
monture est interrogée en premier (\texttt{:Gt\#} et \texttt{:Gg\#}), mais
beaucoup d'installations ne lui ont jamais transmis leur site et elle ne répond
alors rien. Sans site, les éphémérides de la nuit ne peuvent pas être
calculées.

La longitude se saisit \textbf{positive vers l'est}. Le protocole LX200 la
compte à l'envers ; la conversion est faite pour vous quand la monture répond.
\end{warning}

\section{Tolérances et fenêtre glissante}

La tolérance est un seuil d'alerte que vous fixez. Une règle simple : prenez
votre échantillonnage en secondes d'arc par pixel. Une tolérance égale à
l'échantillonnage est stricte, le double est confortable.

La fenêtre de l'écart-type glissant se règle de 60 à 900 secondes. Soixante
secondes réagissent vite, neuf cents lissent une longue session.

\section{L'armement de l'enregistreur}

Deux réglages commandent le comportement du bouton d'enregistrement.

\textbf{Démarrer l'enregistrement quand la monture se met à suivre}, actif par
défaut. Le bouton \emph{arme} alors l'enregistreur au lieu de le démarrer : rien n'est écrit tant
que la monture ne suit pas. Une monture connectée à midi pour une nuit qui
commence à 19~h n'enregistre plus sept heures de rien. Un second clic désarme.

\textbf{Suspendre quand la monture cesse de suivre}, actif par défaut, avec un
délai de grâce de deux minutes. Le délai existe pour une raison : un arrêt peut être
passager --- retournement au méridien, autofocus, slew entre deux cibles ---
et couper au premier échantillon percerait un trou à chaque dither.

\textbf{Terminer la nuit au lever du Soleil}, actif par défaut. Une nuit est un
tout : si la
monture se parque à 2~h et repart à 3~h, l'enregistrement se suspend puis
reprend \emph{dans le même fichier}. Seul le jour clôt la session et déclenche
le rapport --- et seulement si le Soleil est réellement passé sous l'horizon
pendant la session, pour qu'un essai de jour ne se referme pas sur lui-même.

% ─────────────────────────────────────────────────────────────────
\chapter{Une nuit d'observation}
\label{ch:usage}

\section{Déroulement}

\begin{enumerate}
  \item Connecter la monture, vérifier que l'état passe en suivi.
  \item Armer l'enregistreur --- ou le démarrer, si l'armement est désactivé.
  \item Laisser tourner. Le rapport s'écrira tout seul à la fin.
\end{enumerate}

\section{Ce qu'il faut surveiller en direct}

\begin{itemize}
  \item \textbf{L'écart-type glissant}, plus que la déviation instantanée : il
  dit si la monture tient sa position.
  \item \textbf{Les alertes de tolérance}, qui signalent un dépassement.
  \item \textbf{La marge au point de rosée}, si un capteur d'environnement est
  branché.
\end{itemize}

\begin{tip}
Un pic isolé n'est presque jamais la monture : c'est un coup de vent, un câble,
un pas trop près du pied. Ce qui compte est le niveau \emph{entre} les pics.
\end{tip}

\section{Les éphémérides de la nuit}

Calculées depuis le site, elles donnent le coucher, le crépuscule nautique, la
nuit noire, le lever --- et surtout la part de nuit noire que la session a
réellement couverte. Une nuit commencée une heure trop tard perd une heure
qu'on ne rattrape pas, et c'est plus utile de le lire à côté des chiffres de
suivi qu'après coup.

% ─────────────────────────────────────────────────────────────────
\chapter{Les fichiers}
\label{ch:files}

\section{Ce qu'une session écrit}

\begin{center}
\rowcolors{2}{rowalt}{white}
\begin{tabularx}{\textwidth}{lX}
\toprule
\textbf{Extension} & \textbf{Contenu} \\
\midrule
\texttt{.dat} & Les échantillons : heure, AD, DÉC, écarts-types, état. \\
\texttt{.dti} & Les mesures de synchronisation PC / monture. \\
\texttt{.fft} & Les spectres enregistrés en cours de session. \\
\texttt{.env} & Température, pression, humidité, quand la monture les donne. \\
\texttt{.log} & Les événements : connexions, slews, alertes, suspensions. \\
\bottomrule
\end{tabularx}
\end{center}

L'en-tête du \texttt{.dat} porte le site et le décalage horaire. C'est ce qui
permet de recalculer les crépuscules en relecture, longtemps après que la
monture a été débranchée.

\section{Le rapport de nuit}

Il est écrit automatiquement en fin de session, à côté du \texttt{.dat}, en
\textbf{un fichier par langue}. On n'a rien à demander.

\section{La relecture}

\textbf{Fichier $\rightarrow$ Ouvrir un log} relit une session terminée et
reconstruit toute l'analyse : statistiques par cible, mouvements commandés,
dérive, spectres, éphémérides. Les onglets de langue permettent de lire le même
rapport dans les trois langues.

% ─────────────────────────────────────────────────────────────────
\chapter{Quand quelque chose ne va pas}
\label{ch:trouble}

\begin{center}
\rowcolors{2}{rowalt}{white}
\begin{tabularx}{\textwidth}{lX}
\toprule
\textbf{Symptôme} & \textbf{Piste} \\
\midrule
Pas de connexion réseau & Vérifier l'adresse et le port 3492, et que rien
d'autre ne tient déjà la session --- une 10Micron n'en accepte qu'une. \\
Liaison série muette & Les deux bouts doivent s'accorder sur la vitesse. Les
10Micron sortent d'usine en 9600 bauds, et le menu de la monture permet d'en
changer. \\
Dérive en déclinaison & Sur une monture à modèle, ajouter des points près de la
cible. La vis d'azimut n'est pas la réponse. \\
Note médiocre, monture calme & Lire le chapitre~\ref{ch:stats} : c'est souvent
la lecture qui est fausse, pas la monture. \\
Écart-type glissant absurde & Il est normal juste après un slew ou un
repositionnement : sa fenêtre enjambe le mouvement. Le rapport écarte ces
valeurs. \\
Pas d'éphémérides & Le site n'est pas renseigné. Préférences, section site. \\
\bottomrule
\end{tabularx}
\end{center}

% ─────────────────────────────────────────────────────────────────
\chapter{Référence}
\label{ch:ref}

\section{Raccourcis}

\begin{center}
\rowcolors{2}{rowalt}{white}
\begin{tabularx}{\textwidth}{lX}
\toprule
\textbf{Touches} & \textbf{Action} \\
\midrule
\texttt{F1} & Aide \\
\texttt{Ctrl+,} & Préférences \\
\texttt{Ctrl+O} & Ouvrir un log \\
\texttt{Ctrl+Q} & Quitter \\
\bottomrule
\end{tabularx}
\end{center}

\section{Langue}

Le menu \textbf{Langue} propose le français, l'anglais et le néerlandais, plus
la détection automatique depuis le système. Chaque entrée est écrite dans sa
propre langue : qui se retrouve dans la mauvaise ne peut pas lire la langue
courante pour s'en sortir. Le changement prend effet au démarrage suivant.

\section{Signalement des incidents}

Si le programme plante, se fige ou refuse de démarrer, il peut le signaler
lui-même. La question vous est posée une seule fois, et rien d'autre n'est
jamais envoyé : ni vos données de suivi, ni vos noms de fichiers, ni votre nom
d'utilisateur. Tous les chemins sont anonymisés.

\section{Historique des versions}

Le détail des changements, avec les chiffres mesurés qui les justifient, se
trouve dans le fichier \texttt{CHANGELOG.md} livré avec le programme.

\section{Origine}

MountMonitor est une réécriture indépendante en Python et PyQt6 du
MountMonitor v3.37 de \textbf{Nicolàs de Hilster} (2018--2021), écrit en Java.
Aucun code du programme d'origine n'y figure. La licence MIT du présent
programme ne couvre que cette réécriture ; l'\oe uvre originale reste soumise
aux conditions de son auteur.
